\begin{table}
\caption{LASSO-Selektionshäufigkeit je ökonomische Variablengruppe (mittlerer Anteil der Rolling-Fenster, in denen ≥\,1 Variable der Gruppe selektiert wurde). Schock: Energie-Preisschock-Regime (2021-06--2023-03, $n_{\text{Schock}}=22$). Disinflation: 2023-04--2024-10 ($n_{\text{Disinfl.}}=14$). Kostendruck-Hypothese (Cost-Push/Phillips-Kurve): PPI- und LCI-Variablen sollten im Schock-Regime häufiger selektiert werden.}
\label{tab:selection_economic}
\begin{tabular}{lrrr}
\toprule
 & Gesamt & Schock & Disinflation \\
\midrule
PPI (Erzeugerpreise/Cost-Push) & 0.335 & 0.312 & 0.371 \\
BS (Geschäftserwartungen) & 0.200 & 0.178 & 0.234 \\
IP (Industrieproduktion) & 0.200 & 0.201 & 0.199 \\
LCI (Lohnkosten/Cost-Push) & 0.172 & 0.166 & 0.182 \\
ALQ (Arbeitsmarkt/Phillips) & 0.112 & 0.097 & 0.136 \\
\bottomrule
\end{tabular}
\end{table}
